\chapter{2021:Benjamin List and David W.C. MacMillan}

\label{chap:chemistry2021}

The Nobel Prize in Chemistry for the year 2021 was awarded jointly to Benjamin List and David W.C. MacMillan.

\textbf{Motivation:}

for the development of asymmetric organocatalysis

\section{Benjamin List}

\subsection*{Early Life and Education}

Benjamin List was born in 1968 in Frankfurt am Main, Germany.

\subsection*{Career and Major Research}

List received his doctorate from the Goethe University Frankfurt in 1997 and did postdoctoral work with Dieter Seebach at the ETH Zurich. He became director of the Max Planck Institute for Coal Research in Mülheim an der Ruhr. In 2000 he showed that the amino acid proline catalyzes asymmetric aldol reactions, adding a carbon group to a carbonyl compound with high control of stereochemistry.

\subsection*{Nobel Prize-winning Work}

The work for which Benjamin List was awarded the Nobel Prize in Chemistry in 2021 was described by the Nobel Committee as: "for the development of asymmetric organocatalysis".

List's proline catalysis, based on the enamine mechanism used by enzymes, showed that a simple organic molecule could carry out asymmetric synthesis without any metal.

\subsection*{Significance and Impact}

His discovery demonstrated that small organic molecules are powerful asymmetric catalysts, opening the field that is now called organocatalysis.

\subsection*{Later Career and Honors}

List continues his research on organocatalysis and related reactions at the Max Planck Institute for Coal Research in Mülheim.

\subsection*{Legacy}

Organocatalysis has become a standard method for producing enantiomerically pure molecules, used in the synthesis of pharmaceuticals and natural products.

\subsection*{Modern Developments}

List's discovery of asymmetric organocatalysis, showing that small organic molecules can catalyze reactions with high enantioselectivity, founded a new and general field of catalysis. Organocatalysis is now a standard tool in the pharmaceutical and agrochemical industries, enabling the green and cost-effective synthesis of chiral molecules.

\subsection*{Selected Publications}

"Proline-Catalyzed Direct Asymmetric Aldol Reactions" (Journal of the American Chemical Society, 2000)

\clearpage

\section{David W.C. MacMillan}

\subsection*{Early Life and Education}

David W.C. MacMillan was born in 1968 in Bellshill, Scotland.

\subsection*{Career and Major Research}

MacMillan earned a bachelor's degree from the University of Glasgow in 1991 and a doctorate from the University of California, Irvine, in 1996. He became a professor at Princeton University. In 2000 he developed iminium catalysis, in which an amine catalyst activates a carbonyl compound through an iminium ion, and he coined the term organocatalysis for catalysis by small organic molecules.

\subsection*{Nobel Prize-winning Work}

The work for which David W.C. MacMillan was awarded the Nobel Prize in Chemistry in 2021 was described by the Nobel Committee as: "for the development of asymmetric organocatalysis".

MacMillan showed that iminium ions, formed between an amine catalyst and an enal, undergo selective reactions such as cycloadditions with high stereocontrol, and he demonstrated the generality of metal-free asymmetric catalysis.

\subsection*{Significance and Impact}

His methods extended organocatalysis to a broad range of reactions and made it a practical alternative to metal-based asymmetric catalysts.

\subsection*{Later Career and Honors}

MacMillan continues his research on catalysis at Princeton University.

\subsection*{Legacy}

Together with List, MacMillan established organocatalysis as a third pillar of asymmetric synthesis, alongside enzymes and transition-metal catalysts, and his chemistry is now used widely in drug synthesis.

\subsection*{Modern Developments}

MacMillan's development of iminium-enamine organocatalysis, and his coinage of the term organocatalysis, established small-molecule catalysis as a third pillar of catalysis alongside metals and enzymes. Organocatalysis is now widely used in asymmetric synthesis, and his methods enabled efficient routes to chiral drugs and natural products.

\subsection*{Selected Publications}

"New Strategies for Organic Catalysis: The First Highly Enantioselective Organocatalytic Diels-Alder Reaction" (Journal of the American Chemical Society, 2000)

\clearpage

\section*{Chapter Summary}

The 2021 prize recognized the development of asymmetric organocatalysis by Benjamin List and David MacMillan, who showed that small organic molecules can catalyze enantioselective reactions. Organocatalysis is now a major tool for making chiral molecules in the pharmaceutical industry.
